\documentclass[10pt,a4paper]{amsart}
\usepackage[margin=19mm]{geometry}
\usepackage{amsmath,amssymb,mathtools}
\usepackage[scr=rsfs,cal=euler]{mathalfa}
\usepackage{xfrac}
\usepackage[unicode,hidelinks,bookmarks=false]{hyperref}
\newcommand{\ie}{i.e.\ }
\newcommand{\cf}{cf.\ }
\newcommand{\resp}{respectively\ }
\newcommand{\Subsection}{\subsection}
\providecommand{\nobreakdash}{\nobreak\hbox{-}}
\setlength{\emergencystretch}{2em}
\multlinegap0pt
\usepackage{amssymb}
\usepackage{mathtools}
\newtheorem{lemma}{Lemma}
\newcommand{\C}{\mathbb C}
\newcommand{\E}{\mathbb E}
\newcommand{\Prob}{\mathbb P}
\newcommand{\R}{\mathbb R}
\title{Stochastically evolving ellipsoids with symmetries}
\author{Elisha B. Abuya, Nihar Gargava, Yufei Zhao}
\date{Source: arXiv:2606.05105v2 (CC BY 4.0)}
\begin{document}
\maketitle

\noindent\emph{Source and licence.} Stochastically evolving ellipsoids with symmetries. Version arXiv:2606.05105v2 (CC BY 4.0), distributed by arXiv under the Creative Commons Attribution 4.0 International licence, http://creativecommons.org/licenses/by/4.0/, as recorded in the arXiv licence metadata for this exact version and shown on the abstract page https://arxiv.org/abs/2606.05105v2. Full licence notice: \texttt{licenses/arxiv-2606.05105-CC-BY-4.0.txt}.\par\smallskip

\noindent\textit{Editorial scope.} Counted unit: the article's lemma lem:angular-moment (source0.tex lines 1085--1094). The article's complete original proof of it is delivered with it (source0.tex lines 1096--1175, one \texttt{proof} environment of 80 lines). No mathematics is written, repaired or re-derived: the statement and the proof are the article's own lines, in the article's own notation, and no retained line is substituted or re-flowed.  No further source-local context is needed: every cross-reference of every retained line resolves inside the two delivered spans.  Nothing was omitted from any retained span, so the package carries no omission record.  No retained line contains a citation, so the wrapper carries no bibliography and no bibliography entry.  No retained line contains a figure environment, an image-inclusion command or drawing code, so no image is delivered and the package declares no asset dependency.  Editorial formatting only, in addition: an AMS \texttt{amsart} preamble that restates the article's own declarations and macros used by the retained lines, namely the environments \texttt{lemma} on separate counters, together with the standard packages those lines need.  The retained environments print in this document as Lemma 1 in order of appearance, whereas the article prints the same items with the numbering of its own full document.  The complete native source of the article as distributed in the arXiv e-print for this exact version is bundled unchanged as \texttt{sources/202094/source0.tex}.
\begin{lemma}[Angular block-mass moment]\label{lem:angular-moment}
There are universal constants $r_0,c,C>0$ such that, for all $r\ge r_0$ and all
$0\le\lambda\le c\sqrt r$, if
$\omega=(\omega_1,\dots,\omega_s)$ is a uniformly random unit vector in
$S^{N-1}\subset V_\R=\bigoplus_{j=1}^s\C^r$, with $\omega_j\in\C^r$, then
\[
        \E\exp\left(\lambda s\sum_{j=1}^s\|\omega_j\|^4\right)
        \le Ce^\lambda.
\]
\end{lemma}

\begin{proof}
Let $g=(g_1,\dots,g_s)$ be a standard Gaussian vector in
$V_\R=\bigoplus_{j=1}^s\C^r$.  Its direction
$\omega=g/\|g\|$ is uniformly distributed on $S^{N-1}$ and is independent of
$\|g\|$.  If $h$ is a standard Gaussian vector in $\R^d$, then expansion of
$\|h\|^4$ and the identities $\E h_i^2=1$ and $\E h_i^4=3$ give
\[
        \E\|h\|^4=d(d+2).
\]
Since $g_j$ has real dimension $2r$ and $g$ has real dimension $N=2sr$, it
follows that
\[
        \E\|\omega_j\|^4
        =\frac{\E\|g_j\|^4}{\E\|g\|^4}
        =\frac{2r(2r+2)}{N(N+2)}
        =\frac{r(r+1)}{sr(sr+1)}.
\]

Put
\[
        F(\omega)=\left(s\sum_{j=1}^s\|\omega_j\|^4\right)^{1/2}.
\]
By symmetry,
\[
        \E F(\omega)^2
        =s^2\frac{r(r+1)}{sr(sr+1)}
        =\frac{s(r+1)}{sr+1}
        \le1+\frac1r.
\]
Moreover, for $x,y\in S^{N-1}$,
\[
        |F(x)-F(y)|
        \le \sqrt{s}\left(\sum_{j=1}^s
        \bigl|\|x_j\|^2-\|y_j\|^2\bigr|^2\right)^{1/2}
        \le2\sqrt{s}\|x-y\|.
\]
Let $Y=(F(\omega)-\sqrt{1+1/r})_+$.  Since $N=2sr$, concentration on the sphere and
$\E F\le(\E F^2)^{1/2}\le\sqrt{1+1/r}$ imply that there are universal
constants $c,C>0$ such that
\[
        \Prob(Y>u)\le C e^{-2cru^2}
\]
for every $u\ge0$.
Therefore,
\[
        \E e^{crY^2}
        =1+\int_0^\infty 2cru e^{cru^2}\Prob(Y>u)\,du
        \le1+C\int_0^\infty 2cru e^{-cru^2}\,du=O(1).
\]
For $x,y\ge0$, AM--GM gives
\[
        2xy\le\frac{2}{\sqrt r}x^2+\frac{\sqrt r}{2}y^2.
\]
It follows that
\[
        (x+y)^2\le(1+2/\sqrt r)x^2+(1+\sqrt r/2)y^2.
\]
For $0\le\lambda\le c\sqrt r$, applying this with
$x=\sqrt{1+1/r}$ and $y=Y$ gives, after decreasing
$c>0$ if necessary and choosing $r_0$ sufficiently
large,
\[
\begin{aligned}
        \lambda F^2
        &\le\lambda\left(\sqrt{1+1/r}+Y\right)^2 \\
        &\le\lambda(1+2/\sqrt r)(1+1/r)
        +\lambda(1+\sqrt r/2)Y^2 \\
        &\le\lambda+O(1)+crY^2
        && \text{since $\lambda\le c\sqrt r$}.
\end{aligned}
\]
Together with the preceding moment bound, this yields
\[
        \E e^{\lambda s\sum_{j=1}^s\|\omega_j\|^4}
        =\E e^{\lambda F^2}
        \le e^{\lambda+O(1)}\E e^{crY^2}
        =O(e^\lambda).
\]
This proves the lemma.
\end{proof}

\end{document}
